% !TEX TS-program = pdflatex
% !TEX encoding = UTF-8 Unicode

% This file is a template using the "beamer" package to create slides for a talk or presentation
% - Giving a talk on some subject.
% - The talk is between 15min and 45min long.
% - Style is ornate.

% MODIFIED by Jonathan Kew, 2008-07-06
% The header comments and encoding in this file were modified for inclusion with TeXworks.
% The content is otherwise unchanged from the original distributed with the beamer package.

\documentclass{beamer}
\setbeamercovered{invisible}



% Copyright 2004 by Till Tantau <tantau@users.sourceforge.net>.
%
% In principle, this file can be redistributed and/or modified under
% the terms of the GNU Public License, version 2.
%
% However, this file is supposed to be a template to be modified
% for your own needs. For this reason, if you use this file as a
% template and not specifically distribute it as part of a another
% package/program, I grant the extra permission to freely copy and
% modify this file as you see fit and even to delete this copyright
% notice. 


\mode<presentation>
{
  \usetheme{Warsaw}
  % or ...

  % or whatever (possibly just delete it)
}


\usepackage[english]{babel}
% or whatever

\usepackage[utf8]{inputenc}
\usepackage{qrcode}
\usepackage{hyperref}
\usepackage{ytableau}
\usepackage{longtable}
% or whatever

\usepackage{times}
\usepackage[T1]{fontenc}
\usepackage{tikz}
% Or whatever. Note that the encoding and the font should match. If T1
% does not look nice, try deleting the line with the fontenc.

%%% MATH RELATED

\title{MATH 180 Strategies of Problem Solving} 
\subtitle
{} % (optional)

\author[W.R. Casper] % (optional, use only with lots of authors)
{W.R. Casper}
% - Use the \inst{?} command only if the authors have different
%   affiliation.

\institute[California State University Fullerton] % (optional, but mostly needed)
{
  Department of Mathematics\\
  California State University Fullerton}
% - Use the \inst command only if there are several affiliations.
% - Keep it simple, no one is interested in your street address.

\subject{Talks}
% This is only inserted into the PDF information catalog. Can be left
% out. 



% If you have a file called "university-logo-filename.xxx", where xxx
% is a graphic format that can be processed by latex or pdflatex,
% resp., then you can add a logo as follows:

% \pgfdeclareimage[height=0.5cm]{university-logo}{university-logo-filename}
% \logo{\pgfuseimage{university-logo}}



% Delete this, if you do not want the table of contents to pop up at
% the beginning of each subsection:
\AtBeginSubsection[]
{
  \begin{frame}<beamer>{Outline}
    \tableofcontents[currentsection,currentsubsection]
  \end{frame}
}


% If you wish to uncover everything in a step-wise fashion, uncomment
% the following command: 

%\beamerdefaultoverlayspecification{<+->}


\begin{document}

\begin{frame}
  \titlepage
\end{frame}

\begin{frame}{Outline}
  \tableofcontents
  % You might wish to add the option [pausesections]
\end{frame}

% Since this a solution template for a generic talk, very little can
% be said about how it should be structured. However, the talk length
% of between 15min and 45min and the theme suggest that you stick to
% the following rules:  

% - Exactly two or three sections (other than the summary).
% - At *most* three subsections per section.
% - Talk about 30s to 2min per frame. So there should be between about
%   15 and 30 frames, all told.

\section{SoPS Lecture 15}

\begin{frame}{Review}
  \begin{block}{Last time}
    \begin{itemize}
\pause
      \item Stars and bars
\pause
      \item working on Connect 4 problems
    \end{itemize}
  \end{block}
\end{frame}

\subsection{Ready Set Bet}

\begin{frame}{Play the game}
  \begin{block}{Activity}
    \begin{itemize}
\pause
      \item play a couple rounds of Ready Set Bet
\pause
      \item what strategies seem to be the best?
    \end{itemize}
  \end{block}
\end{frame}

\subsection{Expectation Value}

\begin{frame}{Starting example}
  \begin{block}{Problem}
    Your friend proposes the following bet
    \begin{itemize}
\pause
      \item roll two $6$-sided die
\pause
      \item gain a dollar if we roll non-double
\pause
      \item lose $8$ dollars if we roll doubles
    \end{itemize}
    Should you play?
  \end{block}
\end{frame}

\begin{frame}{Example}
  \begin{block}{Problem}
    Your friend proposes the following bet
    \begin{itemize}
\pause
      \item roll two $6$-sided dice
\pause
      \item gain a dollar if we roll non-double
\pause
      \item lose a dollar if we roll doubles
    \end{itemize}
    Should you play?
  \end{block}
\end{frame}

\begin{frame}{Probability Intro}
  \begin{block}{Vocabulary}
    In \textbf{probability theory} we study \textbf{experiments}, like flipping a coin, rolling a die, or drawing a card.\\
    \begin{itemize}
\pause
      \item the set $\Omega$ of all possible outcomes is the \textbf{sample space}
\pause
      \item subsets $E\subseteq \Omega$ are called \textbf{events}
\pause
      \item the \textbf{probability} of an event $E$ is
      $$\mathbb{P}(E) = \frac{|E|}{|\Omega|}$$
    \end{itemize}
  \end{block}
\end{frame}

\begin{frame}{Back to example}
  \begin{block}
\pause
    The \textbf{sample space} is:
    $$\Omega = \{(1,1),(1,2),(1,3),\dots, (6,6)\}$$
\pause
    The \textbf{event} of rolling doubles is:
    $$E = \{(1,1),(2,2),(3,3),(4,4),(5,5),(6,6)\}$$
\pause
    The probability of rolling doubles is
    $$\mathbb{P}(E) = \frac{|E|}{|\Omega|} = \frac{6}{36} = \frac{1}{6}$$
\pause
    The probability of NOT rolling doubles is
    $$\mathbb{P}(E) = 1-\frac{1}{6} = \frac{5}{6}.$$
  \end{block}
\end{frame}

\begin{frame}{Back to example}
  \begin{block}{Solution?}
\pause
    If you take the bet, your probability of winning is about $83\%$.
\pause
    So your very likely to win!  So you should play
  \end{block}
\pause
  $$\text{{\color{red}... or should you?}}$$
\end{frame}

\begin{frame}{Expectation value}
  Generic sample space
  $$\Omega = \{o_1,o_2,o_3,\dots,o_N\}$$
\pause
  \begin{block}{Definition}
\pause
    A \textbf{value function} $f: \Omega\rightarrow\mathbb R$ assigns to each possible experimental outcome, the the value of the outcome.
  \end{block}
\pause
  \begin{block}{Definition}
\pause
    The \textbf{expected value} of a value function $f$ is
    $$\mathbb E(f) = \frac{f(o_1)+f(o_2)+\dots+f(o_N)}{N}.$$
  \end{block}
\pause
  \begin{block}{Meaning}
\pause
    The expected value represents the value of $f$ \emph{on average} if the experiment is conducted lots and lots of times.
  \end{block}
\end{frame}

\begin{frame}{Back to problem}
\pause
  Value function:
  $$f((j,k)) = \left\lbrace\begin{array}{cc}-8, & j=k\\1, & j\neq k\end{array}\right.$$
\pause
  $$f((1,1)) = -8,\ f((2,3)) = 1, f((4,5)) = 1,\dots\text{and so on}$$
\pause
  $$\mathbb E(f) = \frac{\overbrace{1+1 + \dots + 1}^{\text{30 terms}} + \overbrace{(-8) + \dots + (-8)}^{\text{6 terms}}}{36} = \frac{30-48}{36} = -\frac{1}{2}$$
\end{frame}

\begin{frame}{Back to example}
\pause
  \begin{block}{Solution}
\pause
    Each time you take the bet, you lose an average of $50$ cents.\\
\pause
    If you took the bet $100$ times, you should expect to lose around $50$ dollars.
  \end{block}
\pause
  $$\text{{\color{red}Don't take the bet!}}$$
\end{frame}

\begin{frame}{Ready Set Bet}
  \begin{block}{Discussion}
    How can expectation values help us in Ready Set Bet?
  \end{block}
\end{frame}

\subsection{Wrapping up}

\begin{frame}{Final thoughts}
  \begin{block}{Reflection}
    \begin{itemize}
\pause
      \item Revisit the two dice bet where if you roll non-doubles you win a dollar, but if you roll doubles, you lose $m$ dollars.  Which value of $m$ would result in a "fair bet"?  Why?
    \end{itemize}
  \end{block}
\begin{center}
Submit your responses by \textbf{Tonight!}
\end{center}
\end{frame}


\end{document}


